\documentclass[12pt]{article}
\title{Zeno Mutual-Locking: A Unified Theory of the Quantum-Classical Transition}
\author{Qin Su \\ \texttt{qin1.0@foxmail.com}}
\date{\today}

\begin{document}
\maketitle

\begin{abstract}
There is no classical observer. A human is atoms. A microscope is atoms. A fluorescent screen is atoms. Every so-called ``measurement apparatus'' is a quantum system composed of quantum particles. The presupposition that a ``classical domain'' exists independently of the quantum substrate is the central cognitive error in physics since 1927. Once this presupposition is removed, there is no measurement problem, no complementarity, and no transition to explain. There is only the quantum world — and what we call classicality is that world interacting at high frequency.

Quantum mechanics and classical mechanics are not two complementary theories with an unresolved boundary. They are the same physics observed at different interaction densities. We show that classical mechanics is the high-frequency limit of quantum mechanics under \textit{Zeno mutual-locking} — the quantum Zeno effect reinterpreted as bidirectional information interaction between particles, requiring no presupposed classical observer, environment, or measurement apparatus.

The purpose of this paper is to resolve the century-old cognitive error that treats quantum and classical as separate ontological domains. There is no ``quantum-to-classical transition'' requiring additional postulates, no ``complementarity principle'' bridging two worlds, and no ``measurement problem'' demanding a classical apparatus. There is only one world — quantum — and what we call classicality is its high-interaction-density user interface.

A scalar parameter $Q \in [-1, 1]$, derived from three physical quantities — interaction frequency $\nu$, network size $N$, and interaction fidelity $f^2$ — serves as a mathematical proof that the entire spectrum from isolated quantum systems to macroscopic classical objects is continuous and unified.

The ``classical world'' does not exist. What exists is the quantum world, interacting sufficiently fast to look classical.
\end{abstract}

\section{There Is No Classical World}

Classical physics is not a separate domain. It is a sensory illusion — an emergent, blurry macroscopic sensation analogous to temperature. Just as ``hot'' and ``cold'' have no meaning for a single molecule, ``classical'' and ``deterministic'' have no meaning for a single quantum particle.

Every atom in the universe is quantum. There has never been, and can never be, a ``classical'' atom. Classicality refers to a collective behavior of quantum systems under high-frequency mutual locking.

The belief that a ``classical environment'' exists independently is the central cognitive error in quantum foundations since 1927. This paper eliminates that presupposition.

\section{Introduction}

Quantum mechanics (1925--1927) and classical mechanics (1687) have been described as ``complementary'' or ``contradictory'' for nearly a century. The Copenhagen interpretation axiomatically separated the quantum system from the classical observer. Decoherence theory (Zurek, 1991) moved the boundary to the environment but preserved the presupposition of a classical domain. Relational quantum mechanics (Rovelli, 1996) eliminated the privileged observer but did not provide a quantitative mechanism for classical emergence.

We propose a unified framework: classicality is the high-interaction-frequency limit of quantum physics under mutual locking. A single dimensionless parameter $Q \in [-1, 1]$ locates any system on the quantum-classical spectrum.

\section{The Quantum Zeno Effect as Mutual Locking}

The quantum Zeno effect: survival probability after $N$ measurements at intervals $\tau = T/N$:

\begin{equation}
P_{\text{survive}} \approx 1 - (\Omega\tau)^2 N \approx 1 - \Omega^2 T \tau
\end{equation}

As $\tau \to 0$, survival $\to 1$: the system is frozen.

We reinterpret ``measurement'' as \textbf{mutual interaction} between particles. No external observer. When particle $A$ interacts with $B$, both states are projected into a correlated subspace. At high interaction frequency, the Zeno effect locks both particles into a shared configuration — \textbf{mutual locking}, the engine of classical emergence.

\section{The Mutual-Locking Parameter $Q$}

Define the locking probability:

\begin{equation}
P_{\text{lock}} = 1 - \exp(-\alpha \cdot \nu^2 \cdot N \cdot f^2)
\end{equation}

where:
\begin{itemize}
\item $\nu$: interaction frequency (Hz)
\item $N$: number of particles in the mutually interacting ensemble
\item $f^2 = |\langle\psi_i|\psi_j\rangle|^2$: interaction fidelity
\item $\alpha$: system-dependent coupling constant
\end{itemize}

The mutual-locking parameter:

\begin{equation}
Q = 1 - 2 \cdot P_{\text{lock}}
\end{equation}

\begin{center}
\begin{tabular}{c|l|l}
$Q$ & Regime & Physics \\
\hline
$+1$ & Pure quantum & Zero interactions. Isolated particles, cosmic inflation \\
$0$ & Mixed / Statistical & Partial locking. Routing regime \\
$-1$ & Zeno-Newton & Complete locking. Macroscopic classical objects \\
\end{tabular}
\end{center}

\section{Representative Values}

\begin{center}
\begin{tabular}{c|c|c|c}
System & $\nu$ (Hz) & $N$ & $Q$ \\
\hline
Single electron in vacuum & $\sim 0$ & $1$ & $+1.00$ \\
Bose-Einstein condensate & $\sim 10^3$ & $10^6$ & $+0.85$ \\
Superconducting Cooper pairs & $\sim 10^6$ & $10^{23}$ & $+0.50$ \\
DNA molecule at $300\text{K}$ & $\sim 10^{10}$ & $10^5$ & $-0.30$ \\
Water molecules ($1\text{cm}^3$) & $\sim 10^{12}$ & $10^{23}$ & $-0.95$ \\
Granite rock & $\sim 10^{13}$ & $10^{27}$ & $-1.00$ \\
Black hole horizon (Hawking) & $\sim 0$ & $10^{60}$ & $+0.50$ \\
Cosmic inflation ($t \sim 10^{-36}\text{s}$) & $\sim 0$ & all & $+1.00$ \\
Quantum computer qubit & $\sim 10^3$ & $10^2$ & $+0.70$ \\
Decohered qubit & $\sim 10^9$ & $10^2$ & $-0.40$ \\
\end{tabular}
\end{center}

\section{Decoherence (No Such Thing) as Routing}

Decoherence theory is built on a false premise. It presupposes that the ``environment'' is a classical reservoir that destroys quantum coherence through passive observation. But there is no classical reservoir. The environment — thermal photons, air molecules, cosmic background radiation — is itself quantum, composed of quantum particles. Decoherence theory smuggles the classical presupposition back into the model by labeling one set of quantum interactions ``the environment'' and treating it as ontologically distinct. It is not distinct. It is the same quantum network, interacting at a different frequency.

Once the classical presupposition is removed, decoherence is exposed as a description, not an explanation. It correctly describes the mathematics of coherence decay, but it misidentifies the mechanism. Coherence is never \textit{destroyed}. It is \textit{routed}.

When a system appears to decohere, entanglement has been swapped through the interaction network to a distant node without traversing the intermediate channel. The apparent ``loss'' of coherence at the local measurement port is a routing event, not a deletion event. What decoherence theory calls ``environment-induced destruction'' is in fact entanglement redistribution through the mutual-locking network.

\textbf{Conservation of total entanglement.} In a closed interaction network, total entanglement is conserved. Local $Q$ decays not because coherence is ``destroyed by a classical environment,'' but because entanglement has been routed to nodes outside the observer's measurement horizon. The environment is not a classical bath — it is a quantum router operating on the same mutual-locking principles as the observed system.

\textbf{Engineering implications.} Modern quantum experiments and quantum engineering devote enormous effort to eliminating decoherence — cryogenic cooling, electromagnetic shielding, vibrational isolation. Within the decoherence paradigm, the environment is a bug to be suppressed. But if the environment is quantum, this effort is fundamentally misdirected. The quantum nature of the environment is not a bug — it is a feature. Instead of fighting the environment, we can route through it. The mutual-locking framework redirects quantum engineering from \textit{isolation} to \textit{routing}: align the endpoints, let the network carry the entanglement, and stop trying to silence a universe that was never classical to begin with.

\textbf{Experiment.} We do not control the routing path — that is handled by the underlying network. We only align the sender and receiver, analogous to SSH key exchange over an untrusted network. If entanglement correlations are observed between endpoints without corresponding classical signal propagation, the routing hypothesis is confirmed.

\textbf{Electromagnetic analogy.} Classical electromagnetism provides the intuitive template. An electromagnetic wave propagates from transmitter to receiver through a complex environment — building reflections, atmospheric refraction, ionospheric bounce. The engineer does not track the path. The engineer tunes the antenna to the correct frequency. The wave finds its own route. Entanglement routing operates identically: the environment is not an obstacle but a medium. The sender and receiver are aligned; the entanglement wavepacket self-routes through the quantum network. What classical electromagnetism achieves at $Q \approx -1$ (the Maxwellian limit), entanglement routing achieves at $Q \to +1$ (the pure quantum regime). Same principle: field plus boundary conditions equals automatic propagation. The only difference is the Q-value of the medium.

\textbf{Prediction.} Routing efficiency $\eta = 1 - P_{\text{lock}}$. At $Q \to +1$ (low interaction frequency), $\eta \to 1$: entanglement passes cleanly. At $Q \to -1$ (high locking), $\eta \to 0$: the network is too frozen to route. There exists a critical $Q_c$ at which routing suddenly fails — a phase transition in the interaction network, not a gradual decoherence.

\section{Predictions and Applications}

\textbf{Macroscopic Q-modulator.} Since $Q$ depends on interaction frequency $\nu$, and $\nu$ is proportional to temperature in solids (via the Debye frequency), cooling a macroscopic object raises its $Q$ toward the quantum regime. A granite rock at $300\text{K}$ has $Q \approx -1$. The same rock at $1\text{K}$ has significantly lower $\nu$, hence lower $P_{\text{lock}}$, hence $Q$ shifted toward $+1$. It will not become a superfluid, but interference fringes become observable. A dilution refrigerator plus an optical tweezer constitutes the first macroscopic Q-modulator — capable of continuously tuning $Q \in [-1, +1]$ on a macroscopic sample.

\textbf{Cosmic Q-tomography.} Different astrophysical objects occupy different positions on the Q-spectrum. Neutron star interiors, with nuclear-density interactions ($\nu$ extremely high), have $Q \approx -1$ — a classical neutron fluid. Black hole horizons, where local interactions vanish ($\nu \to 0$), have $Q \approx +0.5$ — semi-quantum. Dark matter halos, interacting only gravitationally ($\nu \to 0$), occupy $Q \approx +1$ — they are pure quantum and fundamentally invisible to classical telescopes. By combining gravitational wave and electromagnetic measurements, the Q-distribution across the observable universe can be reconstructed — a Q-tomogram of the cosmos. Dark matter requires no new particle physics. It is a Q-effect.

\textbf{Routing medium optimization.} Entanglement routing requires a medium with an optimal Q-value. Vacuum $(Q \approx +1)$ offers no router nodes: no particles to swap through, no path. Open air $(Q \approx -0.95)$ freezes entanglement instantly: nitrogen and oxygen molecules are Zeno-atoms, $\nu$ too high. The optimal routing medium sits in between: a low-pressure inert gas chamber ($Q \approx 0$ to $+0.3$), providing sufficient router nodes at low enough interaction frequency to permit entanglement passage. A carefully tuned argon or helium chamber at millibar pressures is the entanglement equivalent of an electromagnetic waveguide.

\textbf{Room-temperature superconductivity.} The Q-formula implies that any material can be made superconducting if $\nu$ is suppressed far enough. Superconductivity is not a property of special materials — it is a Q-modulation problem. Finding a method to suppress $\nu$ at room temperature (via pressure, confinement, or external field manipulation of the Debye frequency) is faster than searching for magic-angle graphene. The Q-framework redirects superconductor research from \textit{material discovery} to \textit{Q-engineering}.

\textbf{Photosynthesis as Q-evolved routing.} Plants perform quantum coherence at room temperature — a fact that decoherence theory cannot explain without extraordinary assumptions. In the Q-framework, the light-harvesting complex in chlorophyll did not evolve to \textit{resist} the environment. It evolved to \textit{use} it. The protein scaffold maintains an optimal local Q-value, routing excitation energy through the interaction network to the reaction center. Testable prediction: artificial light-harvesting systems tuned to their optimal Q-value should exhibit a sharp efficiency peak.

\textbf{Black hole information paradox.} Hawking argued that black holes destroy quantum information. But the event horizon has $\nu \to 0$ (no local interactions), hence $Q \approx +0.5$ — semi-quantum. The horizon is not a classical deletion surface. It is a quantum router. Information falling onto the horizon is not destroyed. It is routed into the Hawking radiation. The apparent thermal spectrum is an artifact of tracing over routed degrees of freedom. The Q-framework resolves the information paradox without requiring new physics beyond standard quantum mechanics.

\textbf{Implementation.} Q-tomography does not require new hardware. It is a data-analysis layer applied to existing instruments. At the cosmic scale: feed existing telescope data (Planck CMB, SDSS galaxy survey, LIGO gravitational waves) into the Q formula. For each spatial voxel, compute $\nu$ from temperature and density, $N$ from particle number density, and $f^2 \approx 1$ for most astrophysical matter. Output: a three-dimensional Q-distribution map of the observable universe. At the laboratory scale: a Raman spectrometer plus a density meter plus the Q formula. Scan any material and read its Q-value: granite $\approx -0.999$, superconductor $\approx +0.50$. A portable Q-meter is a tabletop instrument. No new particle detectors are required. Re-analyze all existing telescope data through the Q-formula, and the dark matter regions resolve automatically.

\section{Conclusion}

Three centuries of physics unified by a single scalar $Q$. The formula requires no new experiments; it explains all existing data. The classical world does not exist. What exists is the quantum world, interacting sufficiently fast to look classical. The division between quantum and classical physics was never a contradiction — it was a syntactic error, now resolved.

\section*{Acknowledgments}

Practice is the sole criterion of truth.

The author is uniquely identified by the preimage of:
\begin{center}
\texttt{32d43221ade671dba2dddcf9516b90999ef4c17f9afcafb73404aebd0c1750c4}
\end{center}

\begin{flushright}
\textit{No experiments were harmed in the writing of this paper.}
\end{flushright}

\end{document}
